\section{总结}

本课程设计以三层办公楼为工程对象，建筑面积约\SI{1358.4}{m^2}，平面尺寸约\SI{32.0}{m} $\times$ \SI{14.15}{m}。在参考论文设计思路基础上，本文保留房间级热源参数、冷负荷逐时叠加、末端设备选型和施工图校核的基本流程，同时将建筑扩展为三层，并增加重庆城市分支。计算与绘图均围绕天正暖通设计流程展开：平面图采用黑白CAD出图风格，负荷计算采用CLTD/CLF思路，设备表按房间峰值冷负荷、标准容量档位和楼层组合率生成。

本研究的主要结论如下：

第一，五城市峰值冷负荷为\SI{116.07}{kW}--\SI{150.48}{kW}，冷负荷密度为\SI{85.4}{W/m^2}--\SI{110.8}{W/m^2}。重庆峰值最高，为\SI{150.48}{kW}；深圳为\SI{142.54}{kW}，虽然低于重庆，但供冷季更长；沈阳为\SI{116.07}{kW}，仍需按设计日完成新风和末端冷量校核。

第二，全楼按50个房间设置人员、照明、设备和新风参数，其中47个房间纳入空调服务。设计人数为187人，基准新风量为\SI{7250}{m^3/h}。五城市新风负荷为\SI{46.42}{kW}--\SI{70.55}{kW}，占全楼峰值冷负荷39.6\%--46.9\%。重庆和深圳的新风除湿负荷最突出，DOAS盘管冷量、送风露点和冷凝水排放应作为施工图阶段的重点校核项。

第三，围护结构按城市气候区进行差异化设计。沈阳和天津采用较低的外墙、屋面传热系数；成都、重庆和深圳采用更强的外窗遮阳控制，其中深圳外窗传热系数取\SI{2.20}{W/(m^2\cdot K)}、遮阳系数取0.33。灵敏度分析表明，新风量是影响设备容量的首要参数，30/50~m$^3$/(h$\cdot$人)扰动可使峰值冷负荷变化约9.9\%--11.7\%；窗墙比和设备功率密度次之，围护结构增强和遮阳优化可提供小幅但稳定的容量下降。

第四，VRF+DOAS方案按房间峰值冷负荷乘1.3安全系数选取室内机，三层共配置47--48台室内机；各城市室内机总容量为\SI{200.1}{kW}--\SI{243.1}{kW}，室外机总容量为\SI{174.0}{kW}--\SI{209.5}{kW}，楼层组合率控制在107.7\%--117.9\%。该结果说明房间末端选型、全楼同时峰值和楼层室外机容量不能互相替代，必须分别校核。

综合来看，沈阳适合强调短时峰值和低负荷稳定，天津与成都两类系统均具备可行性，重庆和深圳应优先保证DOAS除湿和长供冷季运行可靠性。本文的核心价值不在于复刻参考论文的具体数值，而在于形成一套可迁移的制冷空调设计流程：黑白CAD平面图确定房间边界，天正暖通房间级参数确定冷负荷，五城市围护结构体现地域差异，设备选型表落实到末端和楼层系统，灵敏度分析用于识别后续施工图最需要复核的设计输入。
